\chapter{Extreme cu legături}

\index{extreme!cu legături}

În seminarul anterior am văzut cum se studiază punctele de extrem
și valorile extreme pentru funcții de 2 sau 3 variabile, definite
pe întreg domeniu, $ \RR^2 $ sau $ \RR^3 $. În multe situații, însă,
domeniul de definiție va fi doar o porțiune a planului sau spațiului,
definită prin (in)ecuații.

Există mai multe posibilități pentru a delimita domeniul de definiție
și le prezentăm în ordinea crescătoare a dificultății.

Considerăm, pentru simplitate, că lucrăm cu o funcție de două variabile:
\[
  f : D \seq \RR^2 \to \RR, \quad f = f(x, y).
\]

\textbf{Pentru toate cazurile cînd domeniul de definiție are o frontieră
  (chiar și interval deschis), studiul se împarte în 2 etape: pe interior
  și pe frontieră.}

Pentru \textbf{domenii de tip dreptunghiular}, $ D = [a, b] \times [c, d] $:
\begin{itemize}
\item se studiază mai întîi pe interiorul domeniului, adică
  $ \dr{Int} D = (a, b) \times (c, d) $, cu metoda din cazul extremelor libere.
  Se impune condiția ca punctele critice să aparțină interiorul domeniului.
\item se studiază apoi pe frontieră ($\dr{Fr}D = \partial D $), care poate fi 
  împărțită în mai multe bucăți: $ x = a, y \in [c, d] $ și celelalte,
  caz pentru care funcția devine de o singură variabilă. În exemplul dat,
  avem $ f(x, y) = f(a, y) = g(y) $, pentru care se studiază extremele
  ca în liceu, fiind vorba de o funcție de o singură variabilă.
\end{itemize}

Pentru \textbf{domenii delimitate de ecuații sau inecuații de gradul întîi}, de
forma $ D = \{ x \in [a, b], m_1 x + n_1 \leq y \leq m_2 x + n_2 \} $, 
studiul se poate face ca mai sus pe interior, iar pe frontieră, se substituie
$ x \in [a, b] $, iar $ y = m_1 x + n_1 $, respectiv $ x \in [a, b] $,
iar $ y = m_2 x + n_2 $, în ambele cazuri avînd funcții de o singură variabilă.

\textbf{Domeniile delimitate de ecuații oarecare} se studiază folosind o metodă
generală (care, de altfel, poate fi folosită și în celelalte cazuri), numită
\emph{metoda multiplicatorilor lui Lagrange}. Astfel, să presupunem că
avem un domeniu de forma $ D : E(x, y) \leq \alpha \in \RR$, cum ar fi, de exemplu,
interiorul unui disc:
\[
  D : x^2 + (y - 2)^2 \leq 9.
\]
\index{extreme!funcția lui Lagrange}
Metoda constă în:
\begin{itemize}
\item studiul pe interiorul domeniului, ca în cazul extremelor libere,
  dar se impune condiția ca punctele critice să satisfacă
  \emph{inegalitatea strictă} $ E(x, y) < \alpha $;
\item pe frontieră, se definește \emph{funcția lui Lagrange}:
  \[
    F(x, y) = f(x, y) - \lambda E(x, y), \quad \lambda \in \RR
  \]
  și se reia calculul, pentru funcția $ F $, cu parametrul $ \lambda $. Mai precis,
  avem de rezolvat sistemul:
  \[
    \begin{cases}
      F_x &= 0 \\
      F_y &= 0 \\
      E(x, y) - \alpha &= 0
    \end{cases} \Longrightarrow (x, y, \lambda), \text{ puncte critice.}
  \]
  Din aceste puncte, extragem perechile $ (x, y) $, pe care le folosim
  mai departe. Rezultatul central este:
\end{itemize}

\begin{theorem}\label{thm:functie-marg}
  O funcție continuă definită pe un domeniu \emph{compact}\footnotemark
  este mărginită și își atinge mar\-gi\-ni\-le.
\end{theorem}
\index{extreme!domeniu compact}

\footnotetext{Compacitatea este un concept matematic destul de subtil și nu
  ușor de înțeles. Intuitiv, însă, va fi suficient să menționăm că domeniile care
  din punct de vedere geometric nu au \qq{găuri} sînt compacte. În particular,
  ceea ce vom folosi în exerciții (discuri, elipse, triunghiuri etc.)
  vor fi toate compacte, deci sîntem în ipotezele teoremei.}

Mai departe, putem continua ca în cazul extremelor libere, i.e.\ cu matricea hessiană
pentru punctele critice $ (x_c, y_c) $ obținute, fie putem cita teorema și va fi
suficient să calculăm $ f(x_c, y_c) $. Aceasta deoarece punctele de extrem
sigur sînt printre punctele critice (teorema lui Fermat din liceu) și deci,
valoarea maximă dintre $ f(x_c, y_c) $ va fi punctul de maxim, iar cea
minimă, punctul de minim.

%%%%%%%%%%%%%%%%%%%%%%%%%%%%%%%%%%%%%%%%%%%%%%%%%%%%%%%%%%%%%%%%%%%%%% 
\section{Exerciții}

\indent\indent 0. Reprezentați grafic domeniile date de următoarele (in)ecuații:
\begin{enumerate}[(a)]
\item $ D_1 = \{ (x, y) \in \RR^2 \mid y \geq x^2, y^2 \leq x \} $
\item $ D_2 = \{ (x, y) \in \RR^2 \mid x \leq 2, 0 \leq y \leq x, xy \leq 1 \} $;
\item $ D_3 = \{ (x, y) \in \RR^2 \mid x^2 + y^2 \leq 2y \} $;
\item $ D_4 = \{ (x, y) \in \RR^2 \mid x^2 + y^2 \leq x, y \geq 0 \} $;
\item $ D_5 = \{ (x, y) \in \RR^2 \mid y \geq 0, x + y - 6 \leq 0, y^2 \leq 8x \} $;
\item $ D_6 = \{ (x, y) \in \RR^2 \mid x^2 + y^2 \leq 2x + 2y - 1 \} $.
\end{enumerate}

\vspace{0.3cm}

1. Fie funcția $ f : D \seq \RR^2 \to \RR, f(x, y) = x^3 + y^3 - 6xy $.
Determinați valorile extreme ale funcției pentru:
\begin{enumerate}[(a)]
\item $ D = [0, 1] \times [0, 1] $;
\item $ D = \{ (x, y) \in \RR^2 \mid x \geq 0, y \geq 0, x + y \geq 5 \} $.
\end{enumerate}

\vspace{0.3cm}

2. Fie funcția $ f : D \to \RR^2, f(x, y) = x^2 + xy + y^2 - 4\ln x - 10\ln y + 3 $.
Determinați valorile extreme ale funcției, dacă $ D $ este domeniul
maxim de definiție (extreme libere!).

\vspace{0.3cm}

3. Fie funcția $ f : D \to \RR^2, f(x, y) = 3xy^2 - x^3 - 15x - 36y + 9 $.
\begin{enumerate}[(a)]
\item Pentru $ D = \RR^2 $, determinați valorile extreme ale funcției.
\item Pentru $ D = [-4, 4] \times [-3, 3] $, determinați valorile
  extreme ale funcției.
\end{enumerate}

\vspace{0.3cm}

4. Fie funcția $ f : D \seq \RR^2 \to \RR, f(x, y) = 4xy - x^4 - y^4 $.
Determinați valorile extreme ale funcției pentru:
\begin{enumerate}[(a)]
\item $ D = \RR^2 $;
\item $ D = [ -1, 2 ] \times [0, 2] $;
\end{enumerate}

\vspace{0.3cm}

5. Fie $ f : D \seq \RR^2 \to \RR, f(x, y) = x^3 + 3x^2y - 15x - 12y $.
Aceeași cerință ca mai sus pentru:
\begin{enumerate}[(a)]
\item $ D = \RR^2 $;
\item $ D = \{ (x, y) \in \RR^2 \mid x \geq 0, y \geq 0, 3y + x \leq 3 \} $.
\end{enumerate}

\vspace{0.3cm}

6. Determinați valorile extreme pentru funcțiile $ f $, definite pe
domeniile $ D $, unde:
\begin{enumerate}[(a)]
\item $ f(x, y) = xy(1 - x - y), \quad D = [0, 1] \times [0, 1] $;
\item $ f(x, y) = x^4 + y^4 - 2x^2 + 4xy - 2y^2, \quad D = (-\infty 0) \times (0, \infty) $;
\item $ f(x, y) = x^3 + 8y^3 - 2xy, \quad D = \{ (x, y) \in \RR^2 \mid x, y \geq 0, y + 2x \leq 2 \} $;
\item $ f(x, y) = x^4 + y^3 - 4x^3 - 3y^2 + 3y, \quad D = \{ (x, y) \in \RR^2 \mid x^2 + y^2 < 4 \} $.
\end{enumerate}

\vspace{0.3cm}

7. (Parțial ETTI, Prof.\ Purtan) Fie funcția: $ f : D \seq \RR^2 \to \RR, f(x, y) = x^2 + 2xy - 4x - y $.
Pentru $ D = \{ (x, y) \in \RR^2 \mid 0 \leq x \leq 1, 0 \leq y \leq -x^2 + 2x \} $,
determinați valorile extreme ale funcției.

\vspace{0.3cm}

8. (Examen ACS, Prof.\ Olteanu) Aflați valorile extreme ale funcției:
\[
  f(x, y) = x^2 + y^2 - 3x - 2y + 1
\]
pe mulțimea $ K : x^2 + y^2 \leq 1 $.

\vspace{0.3cm}

9*. Dintre toate paralelipipedele dreptunghice cu volum constant 1, determinați
pe cel cu aria totală minimă.

\vspace{0.3cm}

10. Fie funcția $ f : D \to \RR, f(x, y, z) = x + 2y - 2z $.
Găsiți extremele funcției pentru $ D = \{ (x, y, z) \in \RR^3 \mid x^2 + y^2 + z^2 \leq 16 \} $.

\vspace{0.3cm}

11. Fie funcția $ f : D \to \RR, f(x, y, z) = -4x - 3y + 6z $.
Găsiți extremele funcției pentru $ D = \{ (x, y, z) \in \RR^3 \mid x^2 + y^2 + z^2 \leq 1 \} $.

\vspace{0.3cm}

12. Fie funcția $ f : D \to \RR, f(x, y) = xy^2(x + y - 2) $. Găsiți extremele funcției
pentru $ D = \{ (x, y) \in \RR^2 \mid x + y \leq 3, x, y \geq 0 \} $.

\vspace{0.3cm}

13. Fie funcția $ f : \RR^2 \to \RR, f(x, y) = xy $. Găsiți extremele funcție pentru
$ D = \{ (x, y) \in \RR \mid x^2 + 2y^2 \leq 1 \} $.

%%% Local Variables:
%%% mode: latex
%%% TeX-master: "../am1-20-21"
%%% End:
